\documentclass[10pt, letterpaper]{article}

% Packages:
\usepackage[
    ignoreheadfoot, % set margins without considering header and footer
    top=2 cm, % seperation between body and page edge from the top
    bottom=2 cm, % seperation between body and page edge from the bottom
    left=2 cm, % seperation between body and page edge from the left
    right=2 cm, % seperation between body and page edge from the right
    footskip=1.0 cm, % seperation between body and footer
]{geometry} % for adjusting page geometry
\usepackage{titlesec} % for customizing section titles
\usepackage{tabularx} % for making tables with fixed width columns
\usepackage{array} % tabularx requires this
\usepackage[dvipsnames]{xcolor} % for coloring text
\definecolor{primaryColor}{RGB}{0, 0, 0} % define primary color
\usepackage{enumitem} % for customizing lists
\usepackage{fontawesome5} % for using icons
\usepackage{amsmath} % for math
\usepackage[
    pdftitle={Rafael Duarte's Master CV},
    pdfauthor={Rafael Duarte},
    pdfcreator={LaTeX},
    colorlinks=true,
    urlcolor=primaryColor
]{hyperref} % for links, metadata and bookmarks
\usepackage[pscoord]{eso-pic} % for floating text on the page
\usepackage{calc} % for calculating lengths
\usepackage{bookmark} % for bookmarks
\usepackage{lastpage} % for getting the total number of pages
\usepackage{changepage} % for one column entries (adjustwidth environment)
\usepackage{paracol} % for two and three column entries
\usepackage{ifthen} % for conditional statements
\usepackage{needspace} % for avoiding page brake right after the section title
\usepackage{iftex} % check if engine is pdflatex, xetex or luatex

% Ensure that generate pdf is machine readable/ATS parsable:
\ifPDFTeX
    \input{glyphtounicode}
    \pdfgentounicode=1
    \usepackage[T1]{fontenc}
    \usepackage[utf8]{inputenc}
    \usepackage{lmodern}
\fi

\usepackage{charter}

% Some settings:
\raggedright
\AtBeginEnvironment{adjustwidth}{\partopsep0pt} % remove space before adjustwidth environment
%\pagestyle{empty} % no header or footer

\usepackage{fancyhdr}
\pagestyle{fancy}
\fancyhf{} % Limpa os cabeçalhos e rodapés padrão
\renewcommand{\headrulewidth}{0.5pt} % Adiciona a linha divisória inferior

% Ajusta o espaço do cabeçalho para caber as 3 linhas sem sobrepor o texto
\setlength{\headheight}{45pt} 
\setlength{\headsep}{15pt} % Espaço entre a linha do cabeçalho e o texto

% Configuração do conteúdo centralizado do cabeçalho
\fancyhead[C]{
    \large \textsc{Curriculum Vitae} \,--\, \textbf{Rafael Mendonça Duarte} \\
    \vspace{3pt}
    \normalsize Information Systems Student at University of São Paulo (USP) \\
    \vspace{3pt}
    \footnotesize \hrefWithoutArrow{https://rmduarte.com.br}{\texttt{rmduarte.com.br}} \hspace{1cm} \hrefWithoutArrow{mailto:rmduarte@usp.br}{\texttt{rmduarte@usp.br}}
}


\setcounter{secnumdepth}{0} % no section numbering
\setlength{\parindent}{0pt} % no indentation
\setlength{\topskip}{0pt} % no top skip
\setlength{\columnsep}{0.15cm} % set column seperation
\pagenumbering{gobble} % no page numbering

\titleformat{\section}{\needspace{4\baselineskip}\bfseries\large}{}{0pt}{}[\vspace{1pt}\titlerule]

\titlespacing{\section}{
    -1pt
}{
    0.3 cm
}{
    0.2 cm
} % section title spacing

\renewcommand\labelitemi{$\vcenter{\hbox{\small$\bullet$}}$} % custom bullet points
\newenvironment{highlights}{
    \begin{itemize}[
        topsep=0.10 cm,
        parsep=0.10 cm,
        partopsep=0pt,
        itemsep=0pt,
        leftmargin=0 cm + 10pt
    ]
}{
    \end{itemize}
} % new environment for highlights

\newenvironment{highlightsforbulletentries}{
    \begin{itemize}[
        topsep=0.10 cm,
        parsep=0.10 cm,
        partopsep=0pt,
        itemsep=0pt,
        leftmargin=10pt
    ]
}{
    \end{itemize}
} % new environment for highlights for bullet entries

\newenvironment{onecolentry}{
    \begin{adjustwidth}{
        0 cm + 0.00001 cm
    }{
        0 cm + 0.00001 cm
    }
}{
    \end{adjustwidth}
} % new environment for one column entries

\newenvironment{twocolentry}[2][]{
    \onecolentry
    \def\secondColumn{#2}
    \setcolumnwidth{\fill, 4.5 cm}
    \begin{paracol}{2}
}{
    \switchcolumn \raggedleft \secondColumn
    \end{paracol}
    \endonecolentry
} % new environment for two column entries

\newenvironment{threecolentry}[3][]{
    \onecolentry
    \def\thirdColumn{#3}
    \setcolumnwidth{, \fill, 4.5 cm}
    \begin{paracol}{3}
    {\raggedright #2} \switchcolumn
}{
    \switchcolumn \raggedleft \thirdColumn
    \end{paracol}
    \endonecolentry
} % new environment for three column entries

\newenvironment{header}{
    \setlength{\topsep}{0pt}\par\kern\topsep\centering\linespread{1.0}
}{
    \par\kern\topsep
} % new environment for the header

% save the original href command in a new command:
\let\hrefWithoutArrow\href

\begin{document}
    \newcommand{\AND}{\unskip
        \cleaders\copy\ANDbox\hskip\wd\ANDbox
        \ignorespaces
    }
    \newsavebox\ANDbox
    \sbox\ANDbox{$|$}
    

    \begin{header}
        \fontsize{25 pt}{25 pt}\selectfont Rafael Mendonça Duarte

        \vspace{5 pt}

        \normalsize
        \mbox{São Carlos, SP, Brazil}%
        \kern 5.0 pt%
        \AND%
        \kern 5.0 pt%
        \mbox{\hrefWithoutArrow{mailto:rmduarte@usp.br}{rmduarte@usp.br}}%
        \kern 5.0 pt%
        \AND%
        \kern 5.0 pt%
        \mbox{\hrefWithoutArrow{tel:+55 (11) 9 4197-0036}{+55 11 9 4197-0036}}%
        \kern 5.0 pt%
        \AND%
        \kern 5.0 pt%
        \mbox{\hrefWithoutArrow{https://rmduarte.com.br}{rmduarte.com.br}}%
        
        \vspace{3 pt}
        
        \mbox{\hrefWithoutArrow{https://linkedin.com/in/rafael-mendonca-duarte}{LinkedIn}}%
        \kern 5.0 pt%
        \AND%
        \kern 5.0 pt%
        \mbox{\hrefWithoutArrow{https://scholar.google.com.br/citations?hl=en&user=PZFZg5EAAAAJ}{Google Scholar}}%
        \kern 5.0 pt%
        \AND%
        \kern 5.0 pt%
        \mbox{\hrefWithoutArrow{http://lattes.cnpq.br/2314499895236009}{Currículo Lattes}}%
        \kern 5.0 pt%
        \AND%
        \kern 5.0 pt%
        \mbox{\hrefWithoutArrow{https://orcid.org/0009-0008-7547-0601}{ORCID}}%
    \end{header}
    \thispagestyle{empty}

    \vspace{5 pt - 0.3 cm}

% ---------------------------------------------------------
% PROFILE SUMMARY
% ---------------------------------------------------------
\section{Profile Summary}
\begin{onecolentry}
Information Systems student at the Institute of Mathematical and Computer Sciences of the University of São Paulo (ICMC-USP) with a strong academic record (GPA 8.6/10, no failures). Currently dedicated to the development and improvement of Graph Convolutional Neural Networks (GCNs) applied to image classification, achieving State-of-the-Art (SOTA) performance in recent research. Experienced in High-Performance Computing, deep learning, software integration, and business process automation.
\end{onecolentry}
\vspace{0.15cm}


% ---------------------------------------------------------
% EDUCATION
% ---------------------------------------------------------
\section{Education}

\begin{twocolentry}{
    Jan 2025 – Dec 2028 (Expected)
}
    \textbf{Bachelor of Science in Information Systems}, University of São Paulo (USP) -- ICMC
\end{twocolentry}

\vspace{0.10 cm}
\begin{onecolentry}
    \footnotesize\textbf{Academic Standing:} GPA: 8.6/10 (No failed courses). \\
    \textbf{Relevant Coursework:} Data Structures, Modern Operating Systems, Object-Oriented Programming (OOP), Calculus, Linear Algebra, Statistics.
\end{onecolentry}

\vspace{0.20 cm}

\begin{twocolentry}{
    2022 – 2024
}
    \textbf{High School}, Anglo Alante Jundiaí
\end{twocolentry}

\vspace{0.15cm}

% ---------------------------------------------------------
% RESEARCH EXPERIENCE
% ---------------------------------------------------------
\section{Research Experience}
        
\begin{twocolentry}{
    Aug 2026 – Present
}
    \textbf{Scientific Initiation Student / Research Intern}, University of São Paulo (USP)
\end{twocolentry}

\vspace{0.10 cm}
\begin{onecolentry}
    \textit{Funded by FAPESP -- São Paulo Research Foundation (Grant \#2026/16921-8). Advisor: Prof. Dr. Lucas Pascotti Valem.}
\begin{highlights}
        \item Conducting research on the project \textit{"Identification of Representative Vertices and Feature Graph Fusion in Graph Convolutional Networks"}.
        \item Investigating graph fusion strategies to effectively combine complementary information from different feature extractors while minimizing noise in Graph Convolutional Networks (GCNs) applied to image classification.
        \item Exploring the insertion and identification of representative (class prototypes) and synthetic (image-cluster-based) vertices to improve graph topology and support model generalization under scarce data or limited labels.
    \end{highlights}
\end{onecolentry}
\vspace{0.10 cm}
\begin{onecolentry}
    \setlength{\leftskip}{0.3cm}
    {\footnotesize \textbf{Key Skills:} \textit{Graph Neural Networks (GNNs), PyTorch, Machine Learning, Data Structures and Data Science.}}
\end{onecolentry}

\vspace{0.15cm}

\begin{twocolentry}{
    Nov 2025 – Jul 2026
}
    \textbf{Scientific Initiation Student / Research Intern}, University of São Paulo (USP)
\end{twocolentry}

\vspace{0.10 cm}
\begin{onecolentry}
    \textit{Funded by "Pró-Reitoria de Pesquisa e Inovação". Advisor: Prof. Dr. Lucas Pascotti Valem.}
\begin{highlights}
        \item Investigated the influence of structural graph properties and message aggregation strategies in GCNs applied to image classification with limited labeled data.
        \item Developed \textbf{GRaNDe} (Gaussian Rank-based Neighborhood Degree), a novel adaptive centrality measure that achieved \textbf{State-of-the-Art (SOTA)} performance, \textbf{improving accuracy by up to 14\%} over existing baselines.
        \item \textbf{First-authored} the paper introducing the proposed method, which was accepted for publication at the Brazilian Symposium on Databases (\textbf{SBBD 2026}).
    \end{highlights}
\end{onecolentry}
\vspace{0.10 cm}
\begin{onecolentry}
    \setlength{\leftskip}{0.3cm}
    {\footnotesize \textbf{Key Skills:} \textit{Graph Neural Networks (GNNs), PyTorch, Machine Learning, Data Structures and Data Science.}}
\end{onecolentry}

\vspace{0.15cm}

% ---------------------------------------------------------
% PROFESSIONAL EXPERIENCE
% ---------------------------------------------------------
\section{Professional Experience}

\begin{twocolentry}{
    Sep 2025 – Oct 2025
}
    \textbf{Technical IT Monitor}, Pró-Aluno Program (USP)
\end{twocolentry}

\vspace{0.10 cm}
\begin{onecolentry}
    \textit{Funded by "Pró-Reitoria de Graduação", "Pró-aluno Program".}
    \begin{highlights}
        \item Provided software and hardware technical support to users in the ICMC undergraduate laboratories.
        \item Assisted students in the use of equipment and ensured the proper functioning and care of the local IT infrastructure.
    \end{highlights}
\end{onecolentry}

\vspace{0.20 cm}

\begin{twocolentry}{
    Jul 2025 – Sep 2025
}
    \textbf{IT Management Intern}, FK Trading
\end{twocolentry}

\vspace{0.10 cm}
\begin{onecolentry}
\begin{highlights}
        \item \textbf{Automation \& Tooling:} Developed custom \textbf{Node.js scripts} for API integrations and web scraping to automate lead generation and client follow-ups. Implemented WhatsApp chatbots and trained 3 out of 5 internal coworkers, replacing manual Kanban funnels and boosting the team's \textbf{operational efficiency by 30\%}.
        \item \textbf{Technical Consulting:} Acted as a consultant for 15 corporate clients, mapping business workflows and standardizing processes to guide their first-time adoption of CRM and ERP systems (RD Station, OMIE).
        \item \textbf{Key Achievement:} Optimized a B2B client's sales pipeline to seamlessly manage up to \textbf{1,000 leads/day} and drive 200 monthly sales. In one critical instance, leveraged automations to help a 5-person sales team clear a backlog of \textbf{800 stalled leads in just 4 days}.
    \end{highlights}
\end{onecolentry}
\vspace{0.10 cm}
\begin{onecolentry}
    \setlength{\leftskip}{0.3cm}
    {\footnotesize \textbf{Key Skills:} \textit{Node.js, API Integration, CRM/ERP (RD Station, OMIE), Business Process Automation, Web Scraping, Mentorship.}}
\end{onecolentry}

\vspace{0.15 cm}

% ---------------------------------------------------------
% EXTENSION & LEADERSHIP PROJECTS
% ---------------------------------------------------------
\section{Extension \& Leadership Projects}

\begin{twocolentry}{
    Jun 2025 – Present
}
    \textbf{Director of High-Performance Computing (HPC) / Coordinator}, Project Egrégora (USP)
\end{twocolentry}

\vspace{0.10 cm}
\begin{onecolentry}
    \textit{IEEE Computer Society Chapter \& EESC - São Carlos School of Engineering.}
    \begin{highlights}
        \item Acting as Student Project Coordinator for the technological extension area aimed at the re-characterization of "TV Box" hardware devices.
        \item Coordinated the assembly of low-cost distributed computing clusters and the implementation of local database servers.
        \item We orchestrated a \textbf{performance and scalability evaluation} of databases on resource-constrained ARM hardware utilizing Docker environments and, we analyzing the latency and throughput trade-offs across different architectures, benchmarking direct DB access with a focus to \textbf{optimize I/O problems}.
    \end{highlights}
\end{onecolentry}
\vspace{0.10 cm}
\begin{onecolentry}
    \setlength{\leftskip}{0.3cm}
    {\footnotesize \textbf{Key Skills:} \textit{Docker, Database Architecture, Performance Benchmarking, System Administration and FastAPI.}}
\end{onecolentry}

\vspace{0.20 cm}

\begin{twocolentry}{
    Mar 2026 – Present
}
    \textbf{Extension Member}, Corvus AI Group (USP)
\end{twocolentry}

\vspace{0.10 cm}
\begin{onecolentry}
    \textit{Project: Combating Digital Disinformation: Fake Image Detection.}
    \begin{highlights}
        \item Acting as a volunteer member researching deep learning techniques in computer vision.
        \item Contributed to the development and continuous optimization of algorithms aimed at detecting fake images and mitigating digital disinformation.
    \end{highlights}
\end{onecolentry}

\begin{twocolentry}{
    Mar 2026 – Present
}
    \textbf{Teaching Coordinator Member}, DATA -- Data Science Study Group (USP)
\end{twocolentry}

\vspace{0.10 cm}
\begin{onecolentry}
    \textit{ICMC -- Institute of Mathematical and Computer Sciences, University of São Paulo.}
    \begin{highlights}
        \item Member of the Teaching Team Coordination of the DATA Group, responsible for planning and delivering introductory Data Science lectures.
        \item Taught \textbf{4 lectures} on \textbf{Introduction to Data Science} and \textbf{scikit-learn} to a class of \textbf{10 students}, applying hands-on analysis over a public Kaggle dataset of \textbf{FIFA World Cup 2026} match data.
    \end{highlights}
\end{onecolentry}
\vspace{0.10 cm}
\begin{onecolentry}
    \setlength{\leftskip}{0.3cm}
    {\footnotesize \textbf{Key Skills:} \textit{scikit-learn, Data Science, Machine Learning, Teaching.}}
\end{onecolentry}

\begin{twocolentry}{
    Feb 2026 – Jun 2026
}
    \textbf{Digital Literacy Instructor for Seniors}, "Informática para a Terceira Idade" (USP)
\end{twocolentry}

\vspace{0.10 cm}
\begin{onecolentry}
    \begin{highlights}
        \item 3rd Edition of the extension activity aimed at the digital inclusion of the external community of São Carlos.
        \item Facilitated classes for the elderly community as a student monitor, teaching essential computer skills, internet navigation, and individually clarifying doubts to promote \textbf{technical autonomy}.
    \end{highlights}
\end{onecolentry}

\vspace{0.15 cm}

% ---------------------------------------------------------
% PUBLICATIONS
% ---------------------------------------------------------
\section{Publications}

\begin{onecolentry}
    \begin{highlightsforbulletentries}
    \item \textbf{DUARTE, R. M.}; PONCIANO, J. R.; VALEM, L. P. \textit{"Gaussian Rank-Based Neighborhood Degree for Graph Neural Networks in Image Classification."} In: \textbf{Proceedings of the 41st Brazilian Symposium on Databases (SBBD 2026)}, São Carlos, SP, Brazil, 2026.
    \begin{itemize}[topsep=0.05cm, parsep=0.05cm, leftmargin=0.5cm, label={--}]
        \item \textbf{Status:} Accepted for publication in Conference Proceedings (Qualis A4).
        \item \textbf{DOI (Preprint):} \url{https://doi.org/10.48550/arXiv.2605.24367}
        \item \textbf{Relevance:} The main contribution is the GRaNDe method, a dynamic centrality measure that integrates neighborhood analysis with Gaussian distance weighting to capture node importance, effectively improving GNNs in image classification.
    \end{itemize}
    
    \vspace{0.15cm}
    
    \item BRITO, G. M.; \textbf{DUARTE, R. M.}; VALEM, L. P. \textit{"Rank Correlation Graphs with Mutual Neighborhood Weighting for Graph Convolutional Networks in Image Classification."} In: \textbf{Pattern Recognition Letters}, 2026.
    \begin{itemize}[topsep=0.05cm, parsep=0.05cm, leftmargin=0.5cm, label={--}]
        \item \textbf{Status:} Published (Qualis A1).
        \item \textbf{DOI:} \url{https://doi.org/10.1016/j.patrec.2026.09.012}
        \item \textbf{Relevance:} Introduces a strategy based on mutual neighborhood for weighting correlation graphs, bringing a comparative analysis of different weighting functions and topologies for consistent image classification.
    \end{itemize}
    \end{highlightsforbulletentries}
\end{onecolentry}

\vspace{0.15 cm}


% ---------------------------------------------------------
% AWARDS & EXTRACURRICULARS
% ---------------------------------------------------------
\section{Awards \& Distinctions}

\begin{onecolentry}
    \begin{highlightsforbulletentries}
        \item \textbf{Academic Excellence Award (2024)} -- Granted by Anglo Alante Jundiaí.
        \item \textbf{Bronze Medal (2024)} -- Kangaroo Math Competition.
        \item \textbf{Bronze Medal (2024)} -- Brazilian Astronomy and Astronautics Olympiad (OBA).
        \item \textbf{Silver Medal (2023)} -- Kangaroo Math Competition.
        \item \textbf{Best Position Paper \& Best Director (2023)} -- High School United Nations Human Rights Council 1 (FACAMP) / FAMUN \& CJONU. Highlighted strong communication, diplomacy, and research skills.
        \item \textbf{Bronze Medal (2022)} -- Brazilian Rocket Exhibition (MOBFOG).
        \item \textbf{Two-time Finalist (4th Phase)} -- National History Olympiad (ONHB).
    \end{highlightsforbulletentries}
\end{onecolentry}


\vspace{0.15 cm}
% ---------------------------------------------------------
% TECHNICAL SKILLS
% ---------------------------------------------------------
\section{Technical Skills}

\begin{onecolentry}
    \textbf{Programming \& Web:} Python, Java, Node.js, SQL, FastAPI. \\
    \textbf{Machine Learning \& Data Science:} PyTorch, Graph Neural Networks (GNNs), NLP (Word2Vec), XGBoost, Pandas. \\
    \textbf{Tools \& Infrastructure:} Docker, Git, Linux, API Integration. \\
    \textbf{Core Concepts:} Data Structures, Object-Oriented Programming (OOP), Database Architecture, Business Process Automation, Statistics.
\end{onecolentry}
\vspace{0.15 cm}
% ---------------------------------------------------------
% LANGUAGES
% ---------------------------------------------------------
\section{Languages}

\begin{onecolentry}
    \textbf{Portuguese:} Native \hspace{0.5cm} | \hspace{0.5cm} \textbf{English:} Intermediate -- proficient in technical reading and writing
\end{onecolentry}
    
\end{document}